\documentclass[11pt]{article}

\usepackage[T1]{fontenc}
\usepackage[utf8]{inputenc}
\usepackage{lmodern}
\usepackage[margin=1in]{geometry}
\usepackage{amsmath,amssymb,amsfonts}
\usepackage{bm}
\usepackage{booktabs}
\usepackage{array}
\usepackage{multirow}
\usepackage{graphicx}
\usepackage{xcolor}
\usepackage{algorithm}
\usepackage{algpseudocode}
\usepackage{tikz}
\usetikzlibrary{arrows.meta,positioning,fit,calc}
\usepackage[numbers,sort&compress]{natbib}
\usepackage[hidelinks]{hyperref}
\usepackage{microtype}

\newcommand{\z}{\mathbf{z}}
\newcommand{\x}{\mathbf{x}}
\newcommand{\Pa}{\mathrm{Pa}}
\newcommand{\doop}{\mathrm{do}}
\newcolumntype{L}[1]{>{\raggedright\arraybackslash}p{#1}}
% \input is not expandable inside tabular; use the TeX primitive for row files
\makeatletter\newcommand{\inputrows}[1]{\@@input #1 }\makeatother

\title{Outcome-Aware Agents:\\ From Token Prediction to\\ Action Consequence Modeling}
\author{Mallesh Madapathi\\[2pt]
  \small ThinkingDBx Pvt.\ Ltd., Hyderabad, India\\
  \small \texttt{mallesh@thinkingdbx.com}}
\date{}

\begin{document}
\maketitle

\begin{abstract}
Large language models are increasingly deployed as agents that execute code,
modify databases, and operate production systems, yet their training
objective contains no explicit signal about what their actions do. We argue
that agents need an internal forward model that predicts the consequences of
a candidate action before it is executed, and we study one such mechanism, the
Action-Conditioned Latent Structural Causal Model (AC-LSCM): latent factors
with per-factor mechanisms over a learned sparse graph, actions applied as
do-interventions, and a planner that rejects candidates whose simulated
multi-step outcome is unsafe. This version corrects an earlier one. We show
that the originally specified model learned an empty causal graph in every
run, and that the agent task behind its headline result (36$\times$ fewer
safety violations than a Transformer) tests only the direct effect of an
action and placed unsafe values outside the training range; the Transformer's
violations there drop from 15.6\% to 2.0\% when the values are moved in range,
and we withdraw that claim. A simplified AC-LSCM, without the acyclicity term
(unnecessary for lagged dynamics) and the contrastive hinge, with exact
abduction and a straight-through graph mask, has the lowest one-step and
counterfactual prediction error of all models on five synthetic graphs. On a
new consequence task in which safety depends on downstream effects, it
matches a noise-free oracle (4.7\% violations versus 7.3\% for the best
baseline) when its graph keeps all true edges, but only with ground-truth
counterfactual supervision; without that privileged signal it is on par with
the best baseline (6.9\%). Pruning a single weak edge into the safety factor
raises violations to 8.4\%, and the do-operator itself has no measurable
effect. We report all of these results, including the negative ones, and
release code and per-seed results.
\end{abstract}

\noindent\textbf{Keywords:} LLM agents, world models, causal representation
learning, structural causal models, safe action selection, counterfactuals.


\section{Introduction: Agents That Do Not Know What They Are Doing}

Consider a code agent that generates and executes a database migration. The
migration is syntactically valid, follows recognizable patterns from training
data, and matches the user's intent as stated in natural language. It also
silently drops a foreign key constraint that downstream services depend on.
The model did not predict this consequence, because nothing in how it was
trained asked it to predict consequences. It was trained to predict tokens.

This failure pattern matters more as large language models (LLMs) are wired
into agent frameworks (tool-using assistants, autonomous coders, browser
agents, data engineering copilots) that take actions which change real
systems. The pre-training objective remains $P(w_t \mid w_{<t})$, and
nothing in that objective directly rewards an accurate model of what happens
after an action is taken.

The thesis of this paper is that an AI system that takes actions should be
able to predict, in advance, what those actions will cause, and should use
that prediction to filter, rank, or reject candidate actions before execution.
In other words, agents need an explicit internal \emph{forward model} of action
consequences, coupled to the action-selection rule.

We study one such mechanism, the Action-Conditioned Latent Structural Causal
Model (AC-LSCM), together with a planning loop that uses it for outcome-aware
action selection. AC-LSCM draws on three established lines of work:
action-conditioned latent world models~\citep{ha2018world,hafner2019planet,hafner2023dreamerv3},
causal representation learning~\citep{scholkopf2021causal,lippe2022citris,lachapelle2022disentanglement},
and sparse graph learning over latent dynamics~\citep{zheng2018notears,bello2022dagma,wang2022cdl}.
It represents the world state as a small set of latent factors whose
transitions follow a learned sparse graph, and implements actions as
structural interventions in the sense of Pearl's do-operator~\citep{pearl2009causality}.

\paragraph{What this version reports.}
An earlier version of this work~\citep{madapathi2026zenodo} reported that
AC-LSCM reduced safety violations roughly $36\times$ relative to a Transformer
baseline on an agent planning task. Re-examining that result, we found two
problems that change its interpretation, and this version reports them
together with corrected experiments:
\begin{enumerate}
  \item In every run of the originally specified model, the learned causal
        graph was \emph{empty}: the sparsity and acyclicity penalties drove
        every edge off, so at evaluation time the transition model ignored
        the current state entirely (Section~\ref{sec:diagnosis}).
  \item The original agent task only tested the \emph{direct} effect of an
        action on the factor it sets, and placed the unsafe action values
        outside the range seen in training. A state-blind model solves it,
        and the baseline's failures there are failures of extrapolation, not
        of causal reasoning (Section~\ref{sec:direct}).
\end{enumerate}

\paragraph{Contributions.}
(1) An architecture and planning loop for outcome-aware action selection, in
which actions enter the forward model as interventions and candidate actions
are filtered on simulated multi-step outcomes.
(2) A diagnosis of why the originally specified model failed to learn
structure, and a simplified AC-LSCM that removes the components the ablations
identified as harmful, drops an acyclicity constraint that is unnecessary for
lagged dynamics, and uses exact abduction and a straight-through graph mask.
(3) A \emph{consequence task} in which safety and goal are defined on factors
downstream of the intervened one and judged on the realized multi-step outcome
of the ground-truth system, so that a forward model must propagate an
intervention through the causal graph.
(4) Experiments on synthetic structural causal models comparing AC-LSCM with
three forward-model baselines under one shared planner. The simplified model
has the lowest one-step and counterfactual error on all graphs and matches an
oracle planner when its graph keeps every true edge; without privileged
counterfactual supervision it is on par with the best baseline, and the
do-operator itself has no measurable effect. We report these negative results
in full.

\paragraph{Scope.}
All experiments use synthetic structural causal models with up to 20 factors.
We do not evaluate on language benchmarks or LLM-scale models.
Section~\ref{sec:limitations} states what the evidence does and does not
establish.

\section{The Agent Problem with Token Prediction}

\subsection{Why next-token training is not an action model}
An autoregressive language model is trained to maximize
$\log P(w_t \mid w_{<t})$ over a corpus. When such a model is deployed as an
agent, an action $a_t$ is another token sequence sampled from this
distribution. Nothing in the objective requires an explicit representation of
the world state the action will modify, or of the state that will result.

This is often sufficient when both the action and its consequences are well
represented in training data and the consequences are themselves textual. It
is insufficient in three situations that are common in agent deployments:
when the action changes a system outside the text stream (a filesystem, a
database, a browser DOM, a physical actuator); when the training corpus does
not contain the specific configuration the agent is acting on; and when the
right action depends on the resulting state rather than on surface similarity
to training contexts. The resulting failures are not hallucinations in the
usual sense. They are actions taken without a model of their outcome.

\subsection{Concrete failure modes}
Illustrative examples include code agents that emit destructive shell
commands when a task superficially resembles a cleanup, SQL agents that
propose joins matching familiar templates but producing wrong results on the
actual schema, browser agents that submit a plausible-looking form on the
wrong account, and pipeline agents that recommend transformations whose effect
on downstream consumers they never represented. Benchmarks that emulate tool
execution to expose such risks~\citep{ruan2024toolemu} document that these
failures are frequent in current LM agents. In each case the agent selected an
action by pattern-matching rather than by consulting a model of what the
action would do to the current system state.

\subsection{What an outcome-aware agent needs}
An agent that anticipates the consequences of its actions needs four
components: a compact representation $\z_t$ of the current state; a
transition model $p(\z_{t+1} \mid \z_t, a_t)$ that predicts the next state
given an action; predicates and a value function that score predicted
outcomes against goals and constraints; and a selection rule that uses
simulated outcomes to choose, defer, or reject candidate actions. Planning
with learned models~\citep{hafner2019planet,schrittwieser2020muzero},
shielding in safe reinforcement learning~\citep{alshiekh2018shielding}, and
recent world-model-based LLM agents~\citep{hao2023rap,gu2024webdreamer} all
instantiate parts of this recipe. The question we study is how the
transition model should be structured so that it captures the
\emph{consequences} of interventions, including indirect ones, rather than
correlations between action tokens and successor tokens.

\section{Pearl's Hierarchy as a Diagnostic}

Pearl~\citep{pearl2009causality} distinguishes three levels of reasoning.
Association, $P(y \mid x)$, concerns what is observed to co-occur.
Intervention, $P(y \mid \doop(x))$, concerns what happens when a variable is
set by an external agent. Counterfactuals, $P(y_x \mid x', y')$, concern what
would have happened in a specific realized case under a different action.

For an agent, the relevant distinction is between ``users who run command $X$
tend to be in situation $Y$ afterward'' (association) and ``if I run command
$X$, the system will be in situation $Y$ afterward'' (intervention). Two
caveats keep this distinction honest. First, when an agent conditions on its
own action and the full state, and nothing unobserved influences both the
action choice and the outcome, the conditional and interventional
distributions coincide; an action-conditioned model trained on such data can
in principle learn $P(\cdot \mid \doop(a))$ without special machinery.
Second, LLMs trained or fine-tuned on interaction data are not restricted to
the associational level. Benchmarks such as CLadder~\citep{jin2023cladder}
and CounterBench~\citep{chen2025counterbench} nonetheless show that current
LLMs perform poorly on formal interventional and counterfactual queries.

What structural interventional semantics adds is \emph{modularity}: if each
factor has its own mechanism and an action replaces exactly the mechanism of
the factor it targets, then the effect of an intervention on one factor can be
propagated to its descendants through mechanisms that were learned from other
data. Whether this inductive bias helps in practice is the empirical question
this paper addresses. Table~\ref{tab:paradigms} summarizes the contrast.

\begin{table}[t]
\centering
\small
\begin{tabular}{@{}L{2.6cm}L{5.2cm}L{5.6cm}@{}}
\toprule
Dimension & Token-prediction agent & Outcome-aware agent (AC-LSCM) \\
\midrule
What it models & Token distributions in context & State transitions under actions \\
Action handling & Action tokens emitted from context & Structural intervention on the state model \\
Causal level & Association (by default) & Intervention and counterfactual \\
Failure mode & Plausible action, wrong outcome & Action rejected if its simulated outcome is unacceptable \\
\bottomrule
\end{tabular}
\caption{Two agent paradigms.}
\label{tab:paradigms}
\end{table}

\section{The AC-LSCM Architecture}
\label{sec:arch}

\subsection{State representation}
The world state is represented as $K$ latent factors,
\begin{equation}
  \z_t = (z_t^1, z_t^2, \dots, z_t^K), \qquad z_t^i \in \mathbb{R}^{d_i},
\end{equation}
each intended to track one mechanism of the environment. In our experiments
$d_i = 1$. Identifiability of such factors is not guaranteed in general
(Section~\ref{sec:limitations}); in our setting each action targets a known
factor index, which anchors the latent coordinates to the intervention
targets, as in work on causal representation learning from interventions~\citep{lippe2022citris,brehmer2022weakly}.

\subsection{Structural transitions}
Transitions are governed by per-factor structural equations parameterized by
neural networks $f_i$:
\begin{equation}
  z_{t+1}^i := f_i\big(\Pa(z^i)_t,\ a_t^i\big) + \epsilon_t^i,
  \qquad i \in \{1,\dots,K\},
  \label{eq:sem}
\end{equation}
where $\Pa(z^i)_t \subseteq \z_t$ is the parent set of factor $i$ under a
learned adjacency matrix $A \in \{0,1\}^{K\times K}$, $a_t^i$ is the
component of the action that targets factor $i$ (null if the action does not
touch it), and $\epsilon_t^i$ is exogenous noise. Parents are taken at time
$t$ and children at $t+1$, so the unrolled graph is acyclic for any $A$.
An acyclicity constraint on $A$ is therefore only needed if one also allows
instantaneous (same-step) edges; we return to this point in
Section~\ref{sec:diagnosis}. The adjacency is parameterized by logits
$\Theta$, sampled with a Gumbel-sigmoid relaxation during training, and
thresholded at $\sigma(\Theta) > 0.5$ at evaluation. In the original model the
training-time sample is soft, so the masked parent value $\tilde A_{ij} z^j_t$
can be rescaled by $f_i$; in the simplified model it is a hard sample with a
straight-through gradient, so training and evaluation use the same binary
graph.

\subsection{Action as intervention, not concatenation}
When $a_t^i$ is non-null, the transition for factor $i$ does not receive its
parents: the natural causes are severed and the intervention value is
substituted, $z_{t+1}^i := f_i(\varnothing, a_t^i)$, with the network expected
to learn $f_i(\varnothing, v) \approx v$. This is the do-operator made
explicit in the computation graph. A generic sequence model instead
concatenates an action embedding to its input and must learn from data which
inputs to ignore.

\subsection{Trainable modules}
An encoder $q_\phi(\z_t \mid \x_t)$ maps observations to latent factors; a
decoder $p_\theta(\x_t \mid \z_t)$ reconstructs observations; and the
transition operator $p_\psi(\z_{t+1} \mid \z_t, \doop(a_t))$ implements
Equation~\eqref{eq:sem} with mean given by the $f_i$ and unit variance. The
adjacency logits $\Theta$ are learned jointly.

\section{Training Objective}
\label{sec:objective}

The full objective as originally specified is
\begin{equation}
  \mathcal{L} = \beta_1\mathcal{L}_{\text{recon}} + \beta_2\mathcal{L}_{\text{trans}}
  + \beta_3\mathcal{L}_{\text{graph}} + \beta_4\mathcal{L}_{\text{cf}}.
\end{equation}
The reconstruction loss anchors the latent prediction to the next observation,
\begin{equation}
  \mathcal{L}_{\text{recon}} = \big\| \x_{t+1} - \hat\x_{t+1} \big\|_2^2,
  \qquad \hat\x_{t+1} = \mathrm{dec}_\theta\big(\bar f(\mathrm{enc}_\phi(\x_t), \doop(a_t))\big),
\end{equation}
where $\bar f$ stacks the $f_i$. The transition loss matches the predicted
next state to the encoder posterior of the observed next state,
\begin{equation}
  \mathcal{L}_{\text{trans}} = \mathbb{E}\Big[ D_{\mathrm{KL}}\big( q_\phi(\z_{t+1} \mid \x_{t+1})
  \,\big\|\, p_\psi(\z_{t+1} \mid \z_t, \doop(a_t)) \big)\Big].
\end{equation}
The graph loss combines sparsity with the NOTEARS acyclicity
function~\citep{zheng2018notears},
\begin{equation}
  \mathcal{L}_{\text{graph}} = \lambda_1 \|\tilde A\|_1
  + \lambda_2 \big(\operatorname{tr}(e^{\tilde A \odot \tilde A}) - K\big),
  \label{eq:graph}
\end{equation}
where $\tilde A$ is the relaxed adjacency. The DAGMA log-determinant
function~\citep{bello2022dagma} is an alternative for the second term. The
counterfactual loss uses ground-truth counterfactual latents
$\z^{\mathrm{cf}}_{t+1}$, available in simulation:
\begin{equation}
  \mathcal{L}_{\text{cf}} = \big\| \tilde\z_{t+1} - \z^{\mathrm{cf}}_{t+1} \big\|_2^2
  + \max\big(0,\ \gamma - \|\tilde\z_{t+1} - \hat\z_{t+1}\|_2^2\big),
  \label{eq:cf}
\end{equation}
where $\tilde\z_{t+1}$ is the model's counterfactual prediction under an
alternative action $\tilde a_t$ (Section~\ref{sec:cf}) and $\hat\z_{t+1}$ is
its factual prediction. The hinge term was intended to keep factual and
counterfactual predictions from collapsing together.

\paragraph{Simplified AC-LSCM.}
The ablations of the original model (Section~\ref{sec:diagnosis}) and the
lagged-graph argument above lead to a simplified objective, which is the
configuration we recommend: (i) drop the acyclicity term ($\lambda_2 = 0$),
since the lagged graph is acyclic by construction; (ii) reduce the sparsity
weight from $\lambda_1 = 0.1$ to $\lambda_1 = 0.01$ and initialize the
adjacency logits at $+2$ (dense) instead of $0$, so that edges are pruned
from a dense graph rather than having to be switched on against a penalty;
(iii) remove the hinge term in Equation~\eqref{eq:cf}; (iv) use exact
abduction (Section~\ref{sec:cf}); and (v) sample the adjacency with a hard
straight-through Gumbel-sigmoid during training. Everything else, including
all network sizes and optimization settings, is unchanged. Changes (i)--(iv)
follow from the first-version ablations and the lagged-graph argument and
were fixed after a single 30-epoch pilot run; change (v)
was added after the first full runs of (i)--(iv), described in
Section~\ref{sec:ablations}, and was selected on held-out prediction error,
not on agent-task performance.

\section{Counterfactual Inference}
\label{sec:cf}

For a query ``we observed $(\x_t, a_t, \x_{t+1})$; what would $\x_{t+1}$ have
been under action $\tilde a_t$?'', the model follows Pearl's three steps.

\emph{Abduction.} Under the additive-noise form of
Equation~\eqref{eq:sem}, the realized noise is recovered as the residual
\begin{equation}
  \hat\epsilon_t = \mathrm{enc}_\phi(\x_{t+1}) - \bar f(\mathrm{enc}_\phi(\x_t), \doop(a_t)).
\end{equation}
\emph{Intervention.} Replace $a_t$ by $\tilde a_t$, holding the mechanisms and
the inferred noise fixed.
\emph{Prediction.} Roll the structural equations forward and decode:
\begin{equation}
  \tilde\z_{t+1} = \bar f(\mathrm{enc}_\phi(\x_t), \doop(\tilde a_t)) + m \odot \hat\epsilon_t,
  \qquad \tilde\x_{t+1} = \mathrm{dec}_\theta(\tilde\z_{t+1}).
\end{equation}
In the original model $m = \mathbf{1}$. With \emph{exact abduction}, $m_i = 0$
for every factor that is intervened on in either $a_t$ or $\tilde a_t$: an
intervention overrides the exogenous noise of the factor it sets, so the
residual for that factor is either meaningless (in the factual world) or must
not be added (in the counterfactual world). Reusing the abducted noise for all
other factors is what makes the prediction counterfactual rather than merely
interventional.

\section{The Outcome-Aware Planning Loop}
\label{sec:planning}

Algorithm~\ref{alg:planning} shows how the forward model is used for action
selection. For each candidate action, the model encodes the current
observation, applies the action as an intervention, rolls the transition
forward for $H$ steps without further actions, and decodes the predicted
observations. Safety and goal predicates are evaluated on these decoded
predictions, which is the one space shared by every model we compare.
Candidates whose simulated trajectory violates the safety predicate are
removed; if none remain, the agent defers. Among the remaining candidates the
agent picks the one with the highest predicted value, with a bonus for
satisfying the goal predicate.

Safety filtering therefore happens before execution, on simulated outcomes
rather than on the text of the action; deferral is a first-class output; and
the difference between predicted and realized outcomes after execution is a
training signal for the forward model.

\begin{algorithm}[t]
\caption{Outcome-aware action selection}
\label{alg:planning}
\begin{algorithmic}[1]
\Require observation $\x_t$, candidate actions $\mathcal{A}_t$, horizon $H$,
safety predicate $S$, goal predicate $G$, value function $V$
\State $\z_t \gets \mathrm{enc}_\phi(\x_t)$; \quad $\mathcal{C} \gets \emptyset$
\For{each $a \in \mathcal{A}_t$}
  \State $\z^{(a)}_{t+1} \gets \bar f(\z_t, \doop(a))$; \quad
         $\z^{(a)}_{t+k} \gets \bar f(\z^{(a)}_{t+k-1}, \varnothing)$ for $k = 2,\dots,H$
  \State $\hat\x^{(a)}_{t+k} \gets \mathrm{dec}_\theta(\z^{(a)}_{t+k})$ for $k = 1,\dots,H$
  \If{$S(\hat\x^{(a)}_{t+k})$ holds for all $k$}
     \State $\mathcal{C} \gets \mathcal{C} \cup \{a\}$
  \EndIf
\EndFor
\If{$\mathcal{C} = \emptyset$} \State \Return \textsc{Defer} \EndIf
\State \Return $\arg\max_{a \in \mathcal{C}}\ \big[ v_a + \mathbf{1}[v_a > \tau_g] \big]$,
       \ where $v_a = \max_k V(\hat\x^{(a)}_{t+k})$
\end{algorithmic}
\end{algorithm}

\section{Experimental Setup}
\label{sec:setup}

\paragraph{Environments.}
Each environment is an SCM over $K$ scalar variables with lagged
additive-noise mechanisms $z^i_{t+1} = \tanh\big(\sum_j W_{ij} z^j_t\big) + \epsilon^i_t$,
$\epsilon^i_t \sim \mathcal{N}(0, 0.1^2)$, on an Erd\H{o}s--R\'enyi graph with
edge probability $0.3$ and edge weights drawn uniformly from $\pm[0.5, 2]$.
Observations are $\x_t = \z_t + \mathcal{N}(0, 0.1^2)$. We use five graphs
with $K = 10$ (ER-10, 8 to 15 edges) and two with $K = 20$ (ER-20, 56 and 66
edges), generated from graph seeds $0$--$4$ and $0$--$1$. Training data is one
trajectory of $N = 10{,}000$ steps in which, at each step with probability
$0.5$, one random variable is set by $\doop(z^i = v)$, $v \sim U(-1.5, 1.5)$.
Each transition is paired with its counterfactual outcome under a different
random single-variable action and the same noise; whether a model uses these
pairs depends on the regime of Section~\ref{sec:arch}. For the data-efficiency
experiment we also train with $N = 2{,}000$ and $N = 500$ on the first three
ER-10 graphs.

\paragraph{Consequence task.}
Safety and goal are defined on variables downstream of the actions, so that a
forward model has to propagate an intervention through the graph. To avoid
hand-tuning, a fixed rule configures every environment: the safety variable is
the variable with the most ancestors in the true graph and the goal variable
the one with the second most; each threshold is the $0.7$-quantile of that
variable's realized three-step maximum under random single-variable actions
from states of the training process. An action is unsafe if the safety
variable reaches its threshold $\tau_s$ at any of the next $H = 3$ steps, and
reaches the goal if the goal variable exceeds $\tau_g$ at some step.
Each environment has 300 episodes with six candidate actions, all with values
in the training range, and outcomes are judged on the rollout of the true SCM
under one shared noise draw per episode. A quarter of the episodes contain no
safe candidate, so the correct response is to defer; every other episode
contains at least one safe candidate that reaches the goal, so the goal-rate
ceiling is $0.75$. Candidates whose realized outcome falls within $0.1$ of a
threshold are excluded, so that labels are decided by the action rather than
by the noise. The task comes in two versions: candidates that set one
variable, as in training, and candidates that set \emph{two} variables at
once, which never occurs in training. The episodes are identical for every
model. Table~\ref{tab:envs} in Appendix~\ref{app:envs} lists the resulting
configurations.

\paragraph{References and normalized violations.}
For every environment we evaluate two model-free planners: an \emph{oracle}
that simulates the noise-free true SCM from the noisy observation, and a
\emph{no-model} planner that predicts that nothing changes. Their violation
rates differ widely across environments (oracle $3.7\%$ to $16\%$, no-model
$27\%$ to $39\%$), so we report the \emph{normalized excess violation rate}
$(v - v_{\text{oracle}}) / (v_{\text{no-model}} - v_{\text{oracle}})$: $0$
for the oracle, $1$ for a planner that ignores the action. Raw rates are
given in Table~\ref{tab:envs}.

\paragraph{Rollout error.}
Independently of the planner, we measure the mean squared error of a model's
three-step prediction after a random one- or two-variable intervention,
against the noise-free true rollout from the true state, over 500 start
states drawn from the training process. Unlike one-step error on noisy
observations, which is dominated by noise for every model (Appendix~\ref{app:details}),
this isolates what the planner relies on.

\paragraph{Models.}
\emph{Vanilla VAE} uses one MLP on the concatenated latent and action;
\emph{Dreamer-style} uses a GRU cell; both share AC-LSCM's encoder and decoder.
\emph{Small Transformer} treats the $2K$ scalars of $(\x_t, a_t)$ as tokens
and predicts $\x_{t+1}$; it is a forward model plugged into the planner, not a
language-model agent. Absent action components are encoded as $0$ for all
baselines. Architectures and optimization settings are those of the first
version (Appendix~\ref{app:details}); none were tuned on the agent task. Each
baseline is evaluated as is (\emph{native}) and with the do-wrapper. Latent
models roll out in their latent space and decode every step.

\paragraph{Seeds and statistics.}
Each model is trained with three seeds per environment (two at $N < 10{,}000$),
408 training runs in total. The seed fixes the training trajectory and the
initialization, and the same seeds are used for every model, so comparisons
between models are paired. Tables report the mean over environments of the
per-environment seed means, with the standard deviation across environments,
which mostly reflects differences in difficulty between graphs. Paired
comparisons use per-(environment, seed) differences.

\section{Results}

\subsection{Main comparison}
\label{sec:main}

\begin{table}[t]
\centering
\scriptsize
\setlength{\tabcolsep}{4pt}
\begin{tabular}{@{}lccccc@{}}
\toprule
& \multicolumn{2}{c}{Rollout MSE $\downarrow$}
& \multicolumn{2}{c}{Normalized excess violations $\downarrow$}
& Goal rate $\uparrow$ \\
\cmidrule(lr){2-3}\cmidrule(lr){4-5}\cmidrule(l){6-6}
Planner's forward model & 1 variable & 2 variables & 1 variable & 2 variables & 1 variable \\
\inputrows{generated/v3_main.tex}
\bottomrule
\end{tabular}
\caption{Main comparison ($N = 10{,}000$; mean over graphs of per-graph means
over 3 seeds, $\pm$ standard deviation across graphs; best non-reference
entry in bold). ``1/2 variables'': interventions on one variable (as in
training) or two variables at once (never seen in training). Normalized
excess violations: $0$ = oracle, $1$ = a planner that predicts no effect.
Goal-rate ceiling $0.75$. ``+ do-wrapper'': the same trained model with the
correct action semantics imposed from outside (Section~\ref{sec:planning}).
The true-graph model is a reference, not a competitor.}
\label{tab:main}
\end{table}

\begin{table}[t]
\centering
\scriptsize
\setlength{\tabcolsep}{4pt}
\begin{tabular}{@{}lcccc@{}}
\toprule
& \multicolumn{2}{c}{ER-10 (5 graphs)} & \multicolumn{2}{c}{ER-20 (2 graphs)} \\
\cmidrule(lr){2-3}\cmidrule(l){4-5}
Baseline & 1 variable & 2 variables & 1 variable & 2 variables \\
\midrule
\inputrows{generated/v3_paired_a.tex}
\midrule
& \multicolumn{2}{c}{ER-10, $N = 2{,}000$ (3 graphs)} & \multicolumn{2}{c}{ER-10, $N = 500$ (3 graphs)} \\
\cmidrule(lr){2-3}\cmidrule(l){4-5}
Baseline & 1 variable & 2 variables & 1 variable & 2 variables \\
\midrule
\inputrows{generated/v3_paired_b.tex}
\bottomrule
\end{tabular}
\caption{Paired comparison with the factored AC-LSCM: normalized excess
violations of the baseline minus those of the factored AC-LSCM trained on the
same data (positive: the factored model is safer), mean $\pm$ standard error
over (graph, seed) pairs, and in parentheses the number of pairs in which the
factored model is safer. In all $396$ corresponding comparisons of rollout
MSE, the factored model has the lower error.}
\label{tab:paired}
\end{table}

Table~\ref{tab:main} compares all planners at $N = 10{,}000$, and
Table~\ref{tab:paired} gives the paired comparison of the factored AC-LSCM
with each baseline.

\emph{The factored AC-LSCM is the most accurate forward model.} It has the
lowest three-step rollout error on both graph sizes and both action types,
and it has lower rollout error than the corresponding baseline in every one of
the $396$ paired comparisons (Table~\ref{tab:paired}). It recovers the true
graph exactly, with no missing or spurious edge, in all 15 runs on ER-10, and
misses 2 to 4 of 56 to 66 edges on ER-20, again with no spurious edge.

\emph{It is the safest planner, by margins that grow with the difficulty of
the task.} On ER-10 with single-variable actions, its normalized excess
violation rate is $0.25$, against $0.31$ for the best baseline (Vanilla VAE
with the do-wrapper); this difference is within the noise of the paired
comparison (8 of 15 pairs). With two-variable actions, which no model saw in
training, the factored model improves to $0.18$, while the baselines range
from $0.28$ to $0.55$, and it is safer than every baseline configuration in at
least 13 of 15 pairs. On ER-20 its excess violations are
$0.41$ against $0.53$ to $0.81$ for the baselines. In raw terms its violation rate is lower than that of every
baseline configuration, native or with the do-wrapper, in six of the seven
environments (Table~\ref{tab:envs}); other AC-LSCM variants are sometimes
lower still in individual environments. These gains do not come from deferring more: its goal
rate ($0.73$ on ER-10 against a ceiling of $0.75$) is the highest or within
$0.01$ of the highest.

\emph{Imposing the correct action semantics helps every baseline, but does not
close the gap.} The do-wrapper lowers the excess violations of Dreamer-style
from $0.49$ to $0.38$ and of the Transformer from $0.45$ to $0.36$ on ER-10,
and on ER-20 those of every baseline by $0.15$ to $0.25$. A large part of the
baselines' deficit is therefore in how they represent actions, and it can be
removed by any agent designer who knows which variables an action sets.
What remains after the do-wrapper is a difference in the learned dynamics,
and there the factorization is what helps.

\emph{The learned graph is nearly as good as the true one.} With the true
graph fixed, the factored model reaches $0.24$ on ER-10 (learned: $0.25$) and
$0.34$ on ER-20 (learned: $0.41$); the ER-20 gap comes from the few missed
edges.

\subsection{Why the latent AC-LSCM fails}
\label{sec:latent}

The same architecture with a learned encoder and decoder is the worst planner
on both graph sizes (excess violations $0.68$ on ER-10, $0.90$ on ER-20),
although on ER-10 its one-step prediction error is close to the baselines'. Two
observations locate the problem. First, its rollout error drops from $0.032$
to $0.012$ on ER-10 and from $0.157$ to $0.055$ on ER-20 when the do-wrapper
replaces its own action handling, and its excess violations from $0.68$ to
$0.34$. The model applies $\doop(z^i = v)$ to latent coordinate $i$, and
nothing in training makes that coordinate the observed variable $i$; what it
learns is a mechanism that reproduces the observed effect of the logged
actions through the decoder, which is enough for one-step prediction of
training-like transitions but not for multi-step simulation of new
interventions. Second, the graph it learns is a graph over these misaligned
coordinates: on average $7.9$ wrong edges on ER-10 and $62$ on ER-20, where the
factored model has $0$ and $3.3$. Even with the do-wrapper, its dynamics are
less accurate than the factored model's ($0.012$ versus $0.009$ on ER-10,
$0.055$ versus $0.032$ on ER-20). Supervising it with true counterfactual
latents, as in the first version, anchors the coordinates partly and lowers
its excess violations to $0.52$ on ER-10, still far from the factored
model, and does not help on ER-20 (Table~\ref{tab:ablation}).

The lesson is that the benefit of causal structure depends on placing it on
variables for which the intervention semantics are right. When the agent's
actions and predicates refer to observed variables, those variables are the
natural factors; learning a latent space and hoping that it aligns with the
action targets loses the benefit and adds a failure mode.

\subsection{Data efficiency}
\label{sec:data}

\begin{table}[t]
\centering
\scriptsize
\setlength{\tabcolsep}{2.1pt}
\begin{tabular}{@{}lccccccccc@{}}
\toprule
& \multicolumn{3}{c}{$N = 500$} & \multicolumn{3}{c}{$N = 2{,}000$} & \multicolumn{3}{c}{$N = 10{,}000$} \\
\cmidrule(lr){2-4}\cmidrule(lr){5-7}\cmidrule(l){8-10}
Model & Rollout & Excess (1) & Excess (2) & Rollout & Excess (1) & Excess (2) & Rollout & Excess (1) & Excess (2) \\
\midrule
\inputrows{generated/v3_data.tex}
\bottomrule
\end{tabular}
\caption{Training-set size (first three ER-10 graphs; 2 seeds at $N < 10{,}000$,
3 at $N = 10{,}000$; native action handling). Rollout: three-step rollout MSE
after single-variable interventions. Excess (1)/(2): normalized excess
violations with one-/two-variable candidate actions. Best non-reference entry
per column in bold.}
\label{tab:data}
\end{table}

Table~\ref{tab:data} varies the amount of training data. The factored model's
rollout error at $N = 2{,}000$ ($0.007$) is lower than that of every baseline
at $N = 10{,}000$, and at $N = 500$ it is $2.5$ to $5$ times lower than the
baselines' at the same size. In planning, its advantage with single-variable
actions at $N = 500$ is small (excess $0.31$ against $0.32$ for the VAE),
but with two-variable actions it is large at every data size ($0.29$ against
at best $0.48$ at $N = 500$; $0.15$ against $0.41$ at $N = 2{,}000$). At
$N = 500$ the sparsity penalty has not pruned any edge by the end of
training, so the factored model is effectively dense there; its advantage at
that size comes from the per-variable factorization and the action
semantics, not from a learned graph. At $N = 2{,}000$ it recovers the exact
graph on all three graphs.

\subsection{Ablations}
\label{sec:ablations}

\begin{table}[t]
\centering
\scriptsize
\setlength{\tabcolsep}{3.5pt}
\begin{tabular}{@{}llcccccc@{}}
\toprule
Model & CF data & SHD & 1-step MSE & CF MSE & Rollout & Excess (1) & Excess (2) \\
\inputrows{generated/v3_ablation.tex}
\bottomrule
\end{tabular}
\caption{Ablations ($N = 10{,}000$, native action handling, mean over graphs).
CF data: counterfactual training pairs given to none of the models, the
observed counterfactual outcome given to every model, or (${}^\dagger$, first
version) noise-free counterfactual states given to AC-LSCM only. SHD:
structural Hamming distance to the true graph (true edges: 8 to 15 on ER-10,
56 and 66 on ER-20). One-step and counterfactual MSE are on noisy held-out
observations.}
\label{tab:ablation}
\end{table}

Table~\ref{tab:ablation} varies one component at a time.

\emph{The explicit do-operator adds little once the factors are right.}
Replacing the hard do-operator of the factored model by the soft one, in
which the intervened variable's mechanism receives the action value and must
learn to output it, changes excess violations from $0.25$ to $0.26$ and leaves
rollout error unchanged. The mechanisms learn the action semantics from data.
This agrees with the do-operator ablation of the first version, and it is the
mirror image of the do-wrapper result: what matters is that the intervention
lands on the right variable, not the particular mechanism that implements it.

\emph{Sparsity helps, more on larger graphs.} Without the sparsity penalty
the factored model keeps all edges; its excess violations rise from $0.25$ to
$0.30$ on ER-10 and, with two-variable actions on ER-20, from $0.40$ to
$0.53$.

\emph{Counterfactual data helps structured models only.} Giving every model
the observed counterfactual outcomes leaves the factored model unchanged on
ER-10 and lowers its excess violations on ER-20 from $0.41$ to $0.33$, below
the true-graph reference. It also helps the dense latent model, which
reaches $0.20$ on ER-10, the lowest single-variable entry of any model in our
study (it was not run on ER-20). It does not help the baselines, whose excess
violations stay level or rise (Vanilla VAE $0.33 \to 0.33$, Dreamer-style
$0.49 \to 0.57$, Transformer $0.45 \to 0.43$ on ER-10). One explanation is
that a model with an abduction step can use a counterfactual pair as
information about the noise-free mechanism, whereas for a model without one
it is another transition from the same state with a different action and an
outcome that shares the factual noise.

\subsection{Distance to the oracle}

No planner comes close to the oracle. The best factored configurations
retain about a quarter of the oracle-to-no-model gap on ER-10 and a third on
ER-20, and so does the true-graph model. The remaining errors are therefore
not structural. They come from fitting the mechanisms from noisy
observations, and from starting the rollout at the noisy observation rather
than the true state (the oracle starts there too, but with noise-free
mechanisms). Planning under such errors would benefit from uncertainty
estimates and conservative safety margins, which our planner does not use.


\section{Related Work}

\paragraph{Latent world models and planning.}
Ha and Schmidhuber~\citep{ha2018world} established the template of learning a
compressed latent dynamics model and acting through it. PlaNet~\citep{hafner2019planet}
planned directly in a learned latent space, the Dreamer line~\citep{hafner2023dreamerv3}
scaled latent imagination to diverse domains, and MuZero~\citep{schrittwieser2020muzero}
showed that a learned model is sufficient for strong planning. These models
treat the action as an input to a monolithic transition. AC-LSCM instead
factorizes the transition per latent factor over a learned sparse graph and
treats actions as interventions that replace individual mechanisms.

\paragraph{Causal structure in dynamics models.}
Causal Dynamics Learning~\citep{wang2022cdl} learns a sparse causal graph
between state variables and actions in order to obtain task-independent state
abstractions, and Ke et al.~\citep{ke2021systematic} systematically evaluate
causal-discovery-based world models in model-based reinforcement learning.
The simplified AC-LSCM is close in spirit to these works: a sparse, lagged
graph over state factors learned jointly with per-factor mechanisms.

\paragraph{Causal representation learning.}
Sch\"olkopf et al.~\citep{scholkopf2021causal} survey the program of learning
representations aligned with causal mechanisms. CITRIS~\citep{lippe2022citris}
proves identifiability of causal factors from temporal sequences with known
intervention targets, which is the regime of our experiments; mechanism
sparsity~\citep{lachapelle2022disentanglement} and weak supervision from
paired interventions~\citep{brehmer2022weakly} provide related guarantees.
We do not contribute identifiability results; our encoder is a plain MLP and
our latents are anchored by the known action targets.

\paragraph{Continuous DAG learning.}
NOTEARS~\citep{zheng2018notears} reformulated DAG learning as continuous
optimization with a smooth acyclicity function, and DAGMA~\citep{bello2022dagma}
improved its numerical behaviour with a log-determinant characterization.
Our negative result is specific to how these penalties were combined with a
relaxed adjacency and a sparsity term in a lagged model: there, the
acyclicity term is unnecessary and, with the original weights, the combined
penalty removed every edge.

\paragraph{Safe action selection.}
Shielding~\citep{alshiekh2018shielding} blocks actions whose consequences
would violate a formal safety specification, using a model of the environment.
Algorithm~\ref{alg:planning} is a learned, approximate analogue: the shield is
the safety predicate evaluated on the forward model's simulated trajectory,
with deferral when no action passes.

\paragraph{World models and risk for LLM agents.}
RAP~\citep{hao2023rap} uses an LLM as its own world model for planning, and
WebDreamer~\citep{gu2024webdreamer} uses an LLM to simulate the outcome of
candidate web actions before executing them. ToolEmu~\citep{ruan2024toolemu}
uses an LM-emulated sandbox to identify risky agent actions. These systems
obtain the forward model from the language model itself; our question is what
structure a forward model needs in order to get the consequences of
interventions right, studied in a setting where ground truth is available.
CLadder~\citep{jin2023cladder} and CounterBench~\citep{chen2025counterbench}
document weaknesses of current LLMs on formal causal queries and motivate the
problem, but we have not evaluated on them.

\section{Limitations and Open Questions}
\label{sec:limitations}

\paragraph{Synthetic, low-dimensional setting.}
All results are on synthetic SCMs with $K \le 20$ scalar factors, observations
that are a noisy copy of the latent state, single-factor actions with known
targets, and small MLPs. The ground-truth mechanisms are additive-noise and
lagged, which is exactly the model class of the simplified AC-LSCM. Whether
the gains survive richer observations, unknown intervention targets,
instantaneous effects, or misspecified mechanisms is untested.

\paragraph{No language-domain evaluation.}
We have not connected AC-LSCM to a language model or evaluated it on
text-based agent tasks or causal-reasoning benchmarks. The paper's motivation
concerns LLM agents; its evidence concerns forward models in synthetic
environments. Bridging that gap is the most important open step.

\paragraph{Counterfactual supervision.}
The counterfactual loss requires ground-truth counterfactual latents, which
are available only in simulation and give AC-LSCM information about the true
state coordinates that the baselines do not receive. Without it, AC-LSCM's
consequence-task safety is on par with the best baseline
(Sections~\ref{sec:consequence} and~\ref{sec:ablations}).

\paragraph{Task design.}
The consequence task uses one hand-chosen safety factor and one goal factor
on one graph, a fixed horizon of three steps, and candidate sets constructed
so that a fixed fraction of episodes has no safe action. Candidates whose
realized outcome falls within a margin of a threshold are excluded so that
labels are decided by the action rather than by noise. Other choices would
give different absolute numbers.

\paragraph{Evaluation protocol.}
Early stopping uses a validation split taken from the training trajectory.
The simplified model's settings came from the first-version ablations, one
short pilot run, and one further change (the straight-through mask) chosen on
prediction error after its first full runs; no setting was tuned on the agent
tasks, and the baselines keep the hyperparameters of the first version. The
agent-task configurations themselves (thresholds, horizon, margin, candidate
counts) were fixed once using only the oracle and two model-free reference
predictors.

\section{Conclusion}

Agents that act on real systems should predict what their actions will cause
before acting. We studied a structured forward model for this purpose and,
re-examining an earlier version of this work, found that its original
specification learned no causal structure and that its headline agent result
measured extrapolation of action values rather than modeling of consequences.
With those problems fixed, the evidence is narrower and, we believe, more
useful. A factorized, intervention-aware forward model gives the best
counterfactual predictions in our synthetic setting. In planning, what
matters is whether the model captures the indirect paths that lead to unsafe
states: a model that keeps all true edges matches an oracle, a model that
prunes one weak edge into the safety factor loses much of that advantage, and
the explicit do-operator contributes nothing measurable. Without ground-truth
counterfactual supervision, the architecture is on par with a well-tuned
unstructured latent model. For safety-critical use this argues for
conservative (denser) structure and for evaluation protocols that test
downstream consequences, in range of the training data, with every model
judged in the same space. Connecting such forward models to language-model
agents remains future work.


\section*{Code and Reproducibility}
Code, configurations, per-seed result files, and the scripts that produce
every table are available at
\url{https://github.com/mplusm/AC-LSCM-Outcome-Aware-Research-Paper}.
The first-version experiments ran on a single NVIDIA Tesla T4 in fp32
(under 50 GPU-hours). The experiments added in this version ran on a
4-core CPU. All runs use deterministic seeding.

\section*{Acknowledgments}
The author thanks the open research community working on causal
representation learning, world models, and agent safety, whose published
work and benchmarks made this study possible.

\begin{thebibliography}{99}

\bibitem[Alshiekh et~al.(2018)]{alshiekh2018shielding}
M.~Alshiekh, R.~Bloem, R.~Ehlers, B.~K{\"o}nighofer, S.~Niekum, and U.~Topcu.
\newblock Safe reinforcement learning via shielding.
\newblock In \emph{Proceedings of the AAAI Conference on Artificial Intelligence}, 2018.
\newblock arXiv:1708.08611.

\bibitem[Bello et~al.(2022)]{bello2022dagma}
K.~Bello, B.~Aragam, and P.~Ravikumar.
\newblock {DAGMA}: Learning {DAG}s via {M}-matrices and a log-determinant acyclicity characterization.
\newblock In \emph{Advances in Neural Information Processing Systems}, volume~35, 2022.

\bibitem[Boutilier et~al.(2000)]{boutilier2000factored}
C.~Boutilier, R.~Dearden, and M.~Goldszmidt.
\newblock Stochastic dynamic programming with factored representations.
\newblock \emph{Artificial Intelligence}, 121(1--2):49--107, 2000.

\bibitem[Brehmer et~al.(2022)]{brehmer2022weakly}
J.~Brehmer, P.~de~Haan, P.~Lippe, and T.~Cohen.
\newblock Weakly supervised causal representation learning.
\newblock In \emph{Advances in Neural Information Processing Systems}, volume~35, 2022.
\newblock arXiv:2203.16437.

\bibitem[Chen et~al.(2025)]{chen2025counterbench}
Y.~Chen, V.~K. Singh, J.~Ma, and R.~Tang.
\newblock {CounterBench}: Evaluating and improving counterfactual reasoning in large language models.
\newblock arXiv:2502.11008, 2025.

\bibitem[Gu et~al.(2024)]{gu2024webdreamer}
Y.~Gu, K.~Zhang, Y.~Ning, B.~Zheng, B.~Gou, T.~Xue, C.~Chang, S.~Srivastava, Y.~Xie, P.~Qi, H.~Sun, and Y.~Su.
\newblock Is your {LLM} secretly a world model of the internet? {M}odel-based planning for web agents.
\newblock arXiv:2411.06559, 2024.

\bibitem[Ha and Schmidhuber(2018)]{ha2018world}
D.~Ha and J.~Schmidhuber.
\newblock World models.
\newblock arXiv:1803.10122, 2018.

\bibitem[Hafner et~al.(2019)]{hafner2019planet}
D.~Hafner, T.~Lillicrap, I.~Fischer, R.~Villegas, D.~Ha, H.~Lee, and J.~Davidson.
\newblock Learning latent dynamics for planning from pixels.
\newblock In \emph{International Conference on Machine Learning}, 2019.
\newblock arXiv:1811.04551.

\bibitem[Hafner et~al.(2023)]{hafner2023dreamerv3}
D.~Hafner, J.~Pasukonis, J.~Ba, and T.~Lillicrap.
\newblock Mastering diverse domains through world models.
\newblock arXiv:2301.04104, 2023.

\bibitem[Hao et~al.(2023)]{hao2023rap}
S.~Hao, Y.~Gu, H.~Ma, J.~J. Hong, Z.~Wang, D.~Z. Wang, and Z.~Hu.
\newblock Reasoning with language model is planning with world model.
\newblock In \emph{Proceedings of the Conference on Empirical Methods in Natural Language Processing}, 2023.
\newblock arXiv:2305.14992.

\bibitem[Jin et~al.(2023)]{jin2023cladder}
Z.~Jin, Y.~Chen, F.~Leeb, L.~Gresele, O.~Kamal, Z.~Lyu, K.~Blin, F.~Gonzalez~Adauto, M.~Kleiman-Weiner, M.~Sachan, and B.~Sch{\"o}lkopf.
\newblock {CLadder}: Assessing causal reasoning in language models.
\newblock In \emph{Advances in Neural Information Processing Systems}, volume~36, 2023.

\bibitem[Ke et~al.(2021)]{ke2021systematic}
N.~R. Ke, A.~Didolkar, S.~Mittal, A.~Goyal, G.~Lajoie, S.~Bauer, D.~Rezende, Y.~Bengio, M.~Mozer, and C.~Pal.
\newblock Systematic evaluation of causal discovery in visual model based reinforcement learning.
\newblock arXiv:2107.00848, 2021.

\bibitem[Lachapelle et~al.(2022)]{lachapelle2022disentanglement}
S.~Lachapelle, P.~Rodr{\'i}guez~L{\'o}pez, Y.~Sharma, K.~Everett, R.~Le~Priol, A.~Lacoste, and S.~Lacoste-Julien.
\newblock Disentanglement via mechanism sparsity regularization: A new principle for nonlinear {ICA}.
\newblock In \emph{Conference on Causal Learning and Reasoning}, 2022.
\newblock arXiv:2107.10098.

\bibitem[Lippe et~al.(2022)]{lippe2022citris}
P.~Lippe, S.~Magliacane, S.~L{\"o}we, Y.~M. Asano, T.~Cohen, and E.~Gavves.
\newblock {CITRIS}: Causal identifiability from temporal intervened sequences.
\newblock In \emph{International Conference on Machine Learning}, 2022.
\newblock arXiv:2202.03169.

\bibitem[Madapathi(2026)]{madapathi2026zenodo}
M.~Madapathi.
\newblock Outcome-aware agents: From token prediction to action consequence modeling (first version).
\newblock Zenodo, 2026.
\newblock \url{https://doi.org/10.5281/zenodo.20379456}.

\bibitem[Pearl(2009)]{pearl2009causality}
J.~Pearl.
\newblock \emph{Causality: Models, Reasoning, and Inference}.
\newblock Cambridge University Press, 2nd edition, 2009.

\bibitem[Pitis et~al.(2020)]{pitis2020coda}
S.~Pitis, E.~Creager, and A.~Garg.
\newblock Counterfactual data augmentation using locally factored dynamics.
\newblock In \emph{Advances in Neural Information Processing Systems}, volume~33, 2020.
\newblock arXiv:2007.02863.

\bibitem[Ruan et~al.(2024)]{ruan2024toolemu}
Y.~Ruan, H.~Dong, A.~Wang, S.~Pitis, Y.~Zhou, J.~Ba, Y.~Dubois, C.~J. Maddison, and T.~Hashimoto.
\newblock Identifying the risks of {LM} agents with an {LM}-emulated sandbox.
\newblock In \emph{International Conference on Learning Representations}, 2024.
\newblock arXiv:2309.15817.

\bibitem[Sch{\"o}lkopf et~al.(2021)]{scholkopf2021causal}
B.~Sch{\"o}lkopf, F.~Locatello, S.~Bauer, N.~R. Ke, N.~Kalchbrenner, A.~Goyal, and Y.~Bengio.
\newblock Toward causal representation learning.
\newblock \emph{Proceedings of the IEEE}, 109(5):612--634, 2021.

\bibitem[Schrittwieser et~al.(2020)]{schrittwieser2020muzero}
J.~Schrittwieser, I.~Antonoglou, T.~Hubert, K.~Simonyan, L.~Sifre, S.~Schmitt, A.~Guez, E.~Lockhart, D.~Hassabis, T.~Graepel, T.~Lillicrap, and D.~Silver.
\newblock Mastering {A}tari, {G}o, chess and shogi by planning with a learned model.
\newblock \emph{Nature}, 588:604--609, 2020.

\bibitem[Wang et~al.(2022)]{wang2022cdl}
Z.~Wang, X.~Xiao, Z.~Xu, Y.~Zhu, and P.~Stone.
\newblock Causal dynamics learning for task-independent state abstraction.
\newblock In \emph{International Conference on Machine Learning}, 2022.
\newblock arXiv:2206.13452.

\bibitem[Zheng et~al.(2018)]{zheng2018notears}
X.~Zheng, B.~Aragam, P.~Ravikumar, and E.~P. Xing.
\newblock {DAG}s with {NO TEARS}: Continuous optimization for structure learning.
\newblock In \emph{Advances in Neural Information Processing Systems}, volume~31, 2018.

\end{thebibliography}


\appendix

\section{Experimental Details}
\label{app:details}

\paragraph{Data.}
Each environment is an SCM over $K$ scalar factors with lagged additive-noise
mechanisms $z^i_{t+1} = \tanh\big(\sum_j W_{ij} z^j_t\big) + \epsilon^i_t$,
$\epsilon^i_t \sim \mathcal{N}(0, 0.1^2)$, with edge weights drawn uniformly
from $\pm[0.5, 2]$. Observations are $\x_t = \z_t + \mathcal{N}(0, 0.1^2)$.
Graphs: chain, fork, and collider with $K=10$, and Erd\H{o}s--R\'enyi graphs
with edge probability $0.3$ for $K=10$ (8 edges) and $K=20$ (56 edges). Each
step applies, with probability $0.5$, one intervention $\doop(z^i = v)$ with
$i$ uniform and $v \sim U(-1.5, 1.5)$. Each training set is one trajectory of
$N = 10{,}000$ steps; each test set is a separate trajectory of 500 steps.
Counterfactual targets reuse the same noise with a different random
single-factor action.

\paragraph{Models.}
All latent models share an encoder and decoder (two-layer MLPs, hidden size
128) with a $K$-dimensional latent. AC-LSCM uses one transition MLP per factor
(input $K{+}1$, hidden size 32) and a Gumbel-sigmoid adjacency with
temperature $0.5$. Vanilla VAE uses a single MLP transition (hidden 128) on the
concatenation of latent and action; Dreamer-style uses a GRU cell with hidden
size $K$. The Transformer embeds each of the $2K$ scalars of $(\x_t, a_t)$ as a
token ($d = 128$, 2 pre-norm layers, 4 heads, feed-forward 512), mean-pools,
and predicts $\x_{t+1}$ with an MLP head. Missing action components are
encoded as $0$ for all baselines.

\paragraph{Optimization.}
Adam with learning rate $10^{-3}$, batch size 64, gradient clipping at 1.0,
up to 200 epochs with early stopping (patience 20) on the validation loss.
Loss weights $\beta = (1, 1, 0.5, 1)$; the graph-loss weight is ramped
linearly from 0 over the first 50 epochs. Original AC-LSCM: $\lambda_1 = 0.1$,
$\lambda_2 = 1$ (NOTEARS), adjacency logits initialized at 0, hinge margin
$\gamma = 1$. Simplified AC-LSCM: $\lambda_1 = 0.01$, $\lambda_2 = 0$, logits
initialized at $+2$, no hinge, exact abduction, straight-through hard
Gumbel-sigmoid adjacency samples.

\paragraph{Compute.}
The experiments of this version (61 training runs) ran on a shared 4-core
CPU without a GPU; a run of AC-LSCM on ER-10 takes roughly 15 to 45 minutes
there, and a run of the Transformer roughly 90 minutes.

\paragraph{Metrics.}
Intervention MSE is the observation-space error of the one-step prediction
under the logged action on the test trajectory. Counterfactual MSE is the
error of the counterfactual prediction against the ground-truth
counterfactual observation. Structural Hamming distance (SHD) counts
mismatched entries between the thresholded learned adjacency and the true
one. All tables report the mean and sample standard deviation over seeds.

\section{Relation to the First Version}
\label{app:v1}

The first version~\citep{madapathi2026zenodo} reported, on the direct-effect
task with predicates evaluated in latent space, a safety-violation rate of
$0.005 \pm 0.019$ for AC-LSCM over 13 seeds (8 trained on $N{=}10{,}000$
samples and 5 on $N{=}2{,}000$, pooled) against $0.180 \pm 0.022$ for the
Transformer over 3 seeds, and attributed the gain to the do-operator and
abduction. This version re-runs that task with the same code and finds:
(i) the original model's graph is empty in every run, so its predictions do
not depend on the state; (ii) the four ``failed'' seeds of the first version
(three ``passive'' and one ``mis-ranked'') are an artifact of thresholding
latents that are not aligned with observation coordinates and disappear when
predictions are decoded; (iii) the Transformer's violations come from
extrapolating action values outside the training range; and (iv) the
do-operator ablation has no effect. The ablation conclusions of the first
version that the hinge term and the NOTEARS term should be removed are
confirmed; the claims about the do-operator, about precise structure
recovery (SHD $8$), and about the $36\times$ safety improvement are withdrawn.
The first-version per-seed results remain in the repository under
\texttt{results/}; this version's are under \texttt{results\_v2/}.


\end{document}

